\section{Estimands and Identification}
\label{app:estimands}

This appendix fixes the formal objects used by CVaR-CJE. It contains only
estimand definitions, identification claims, and implementation invariants
that follow from the math.

\subsection{Probability Laws and Mean-CJE}

Let $P_X$ be the prompt law. A policy $\ptarget$ maps prompts to answers:
\[
\X\sim P_X,\qquad
\A_\ptarget\sim\ptarget(\cdot\mid \X),\qquad
\Sscore=s(\X,\A_\ptarget),\qquad
\Y=y(\X,\A_\ptarget).
\]
The cheap judge score is $\Sscore$ and the oracle outcome is $\Y$. Write
$P_\ptarget$ for the joint law of $(\X,\A_\ptarget,\Sscore,\Y)$ and
$\E_\ptarget$ for expectation under that law. Let
\[
\mathcal{G}=\sigma(\Sscore,\X)
\]
be the information available to the calibrator. Under the reference policy
$\pzero$, define
\begin{align}
f_0(\Sscore,\X)
&=\E_{\pzero}[\Y\mid\mathcal{G}],\\
h_{0,t}(\Sscore,\X)
&=\E_{\pzero}\!\left[\shortfall{t}{\Y}\mid\mathcal{G}\right].
\end{align}
When the reference law is clear, write $f=f_0$ and $h_t=h_{0,t}$.

The mean value of policy $\ptarget$ is
\begin{equation}
V(\ptarget)=\E_\ptarget[\Y].
\end{equation}
For any fitted surrogate $f(\Sscore,\X)$,
\begin{equation}
V(\ptarget)-\E_\ptarget[f(\Sscore,\X)]
=\E_\ptarget[\Y-f(\Sscore,\X)].
\label{eq:mean-transport-identity}
\end{equation}

\begin{assumption}[Mean transport]
\label{ass:mean-transport}
For target policy $\ptarget$,
\[
\E_\ptarget[\Y-f(\Sscore,\X)]=0.
\]
\end{assumption}

\begin{proposition}[Mean-CJE identification]
\label{prop:formal-mean-identification}
For any fixed surrogate $f$, \cref{ass:mean-transport} holds if and only if
$\E_\ptarget[f(\Sscore,\X)]=V(\ptarget)$.
\end{proposition}

\begin{twocolproof}
\twocolrow{
\[
V(\ptarget)-\E_\ptarget[f(\Sscore,\X)]
=\E_\ptarget[\Y]-\E_\ptarget[f(\Sscore,\X)].
\]
}{
Use the definitions of $V(\ptarget)$ and the surrogate value.
}
\twocolrow{
\[
\E_\ptarget[\Y]-\E_\ptarget[f(\Sscore,\X)]
=\E_\ptarget[\Y-f(\Sscore,\X)].
\]
}{
Apply linearity of expectation.
}
\twocolrow{
\[
\E_\ptarget[f(\Sscore,\X)]=V(\ptarget)
\Longleftrightarrow
\E_\ptarget[\Y-f(\Sscore,\X)]=0.
\]
}{
This is exactly \cref{ass:mean-transport}.
}
\end{twocolproof}

\subsection{Lower-Tail CVaR}

Fix $\alpha\in(0,1]$. Let
\[
F_\ptarget(t)=\Pp_\ptarget(\Y\le t),
\qquad
q_\alpha(\ptarget)=\inf\{t:F_\ptarget(t)\ge\alpha\}.
\]
The lower-tail Rockafellar--Uryasev objective is
\begin{equation}
\Psi_{\alpha,\ptarget}(t)
=t-\frac{1}{\alpha}\E_\ptarget[\shortfall{t}{\Y}],
\qquad
\shortfall{t}{\Y}=(t-\Y)_+.
\label{eq:app-stoploss-cvar}
\end{equation}
The lower-tail CVaR is
\begin{equation}
\CVaR_\alpha(\ptarget)=\sup_{t\in\R}\Psi_{\alpha,\ptarget}(t).
\label{eq:app-cvar-definition}
\end{equation}

The positive part in $(t-\Y)_+$ is essential. Without it,
\[
t-\E[t-\Y]=t-(t-\E[\Y])=\E[\Y],
\]
so the threshold would carry no tail information and the objective would
collapse to the mean for every $\alpha$.

\begin{assumption}[Regular lower tail]
\label{ass:regular-tail}
For $\alpha\in(0,1)$, the target distribution of $\Y$ is continuous at
$q_\alpha(\ptarget)$, has positive density near $q_\alpha(\ptarget)$, and the
maximizer of \cref{eq:app-stoploss-cvar} is unique up to quantile equivalence.
\end{assumption}

\begin{proposition}[Stop-loss derivation]
\label{prop:formal-stoploss}
Under \cref{ass:regular-tail}, the optimizer of
\cref{eq:app-stoploss-cvar} is $t_\alpha^\star=q_\alpha(\ptarget)$, and
\[
\CVaR_\alpha(\ptarget)
=\E_\ptarget[\Y\mid \Y\le q_\alpha(\ptarget)].
\]
\end{proposition}

\begin{twocolproof}
\twocolrow{
\[
\frac{d}{dt}\E_\ptarget[\shortfall{t}{\Y}]
=\Pp_\ptarget(\Y\le t)=F_\ptarget(t).
\]
}{
At continuity points, differentiating $(t-\Y)_+$ in $t$ gives
$\one\{\Y\le t\}$.
}
\twocolrow{
\[
\frac{d}{dt}\Psi_{\alpha,\ptarget}(t)
=1-\frac{F_\ptarget(t)}{\alpha}.
\]
}{
Differentiate the stop-loss objective.
}
\twocolrow{
$\Psi_{\alpha,\ptarget}$ increases while $F_\ptarget(t)<\alpha$ and decreases
after $F_\ptarget(t)>\alpha$.
}{
The derivative changes sign at the $\alpha$-quantile.
}
\twocolrow{
\[
t_\alpha^\star=q_\alpha(\ptarget).
\]
}{
\cref{ass:regular-tail} gives a regular unique optimizer.
}
\twocolrow{
\[
\begin{aligned}
\Psi_{\alpha,\ptarget}(q_\alpha)
&=q_\alpha-\frac{1}{\alpha}
\E_\ptarget[(q_\alpha-\Y)\one\{\Y\le q_\alpha\}]\\
&=\frac{1}{\alpha}
\E_\ptarget[\Y\one\{\Y\le q_\alpha\}].
\end{aligned}
\]
}{
Continuity gives $\Pp_\ptarget(\Y\le q_\alpha)=\alpha$, so the threshold
terms cancel.
}
\twocolrow{
\[
\Psi_{\alpha,\ptarget}(q_\alpha)
=\E_\ptarget[\Y\mid \Y\le q_\alpha].
\]
}{
Divide by the lower-tail probability $\alpha$.
}
\end{twocolproof}

\subsection{Direct CVaR-CJE Identification}

For each $t$, CVaR estimation is a mean estimation problem for the shortfall
variable $Z_t=\shortfall{t}{\Y}$. The reference-law shortfall surrogate is
\[
h_t(\Sscore,\X)=\E_{\pzero}[Z_t\mid\mathcal{G}].
\]
For target policy $\ptarget$, define the surrogate stop-loss objective
\begin{equation}
\Psi^h_{\alpha,\ptarget}(t)
=t-\frac{1}{\alpha}\E_\ptarget[h_t(\Sscore,\X)].
\label{eq:formal-surrogate-stoploss}
\end{equation}

\begin{assumption}[Shortfall transport]
\label{ass:shortfall-transport}
On an index set $\mathcal{T}\subseteq\R$,
\[
\E_\ptarget[\shortfall{t}{\Y}-h_t(\Sscore,\X)]=0
\qquad\text{for every }t\in\mathcal{T}.
\]
\end{assumption}

\begin{theorem}[CVaR-CJE identification]
\label{thm:formal-cvar-identification}
If \cref{ass:shortfall-transport} holds on $\mathcal{T}$, then
\[
\Psi^h_{\alpha,\ptarget}(t)=\Psi_{\alpha,\ptarget}(t)
\qquad\text{for every }t\in\mathcal{T}.
\]
Consequently,
\[
\sup_{t\in\mathcal{T}}\Psi^h_{\alpha,\ptarget}(t)
=\sup_{t\in\mathcal{T}}\Psi_{\alpha,\ptarget}(t).
\]
If $\mathcal{T}=\R$, the Direct CVaR-CJE population target equals
$\CVaR_\alpha(\ptarget)$.
\end{theorem}

\begin{twocolproof}
\twocolrow{
\[
\Psi_{\alpha,\ptarget}(t)-\Psi^h_{\alpha,\ptarget}(t)
=-\frac{1}{\alpha}
\E_\ptarget[\shortfall{t}{\Y}-h_t(\Sscore,\X)].
\]
}{
Subtract the surrogate objective from the oracle objective; the threshold
terms cancel.
}
\twocolrow{
\[
\Psi_{\alpha,\ptarget}(t)-\Psi^h_{\alpha,\ptarget}(t)=0.
\]
}{
Apply \cref{ass:shortfall-transport} at the same threshold $t$.
}
\twocolrow{
\[
\sup_{t\in\mathcal{T}}\Psi^h_{\alpha,\ptarget}(t)
=\sup_{t\in\mathcal{T}}\Psi_{\alpha,\ptarget}(t).
\]
}{
Identical functions on the same index set have identical suprema.
}
\end{twocolproof}

\begin{assumption}[Nuisance and grid regularity]
\label{ass:nuisance}
The fitted shortfall calibrators are independent of the evaluation folds,
uniformly square-integrable near the relevant maximizers, and satisfy
\[
\sup_{t\in\mathcal{T}_n}
\left|
\frac{1}{n}\sum_{i=1}^n \hat h_t(\Sscore_i,\X_i)
-\E_\ptarget[h_t(\Sscore,\X)]
\right|=o_p(1).
\]
The grid $\mathcal{T}_n$ is fixed before policy comparisons and covers the
population maximizers at the target precision.
\end{assumption}

\subsection{Discrete Outcomes and Atom Splitting}
\label{sec:app-discrete}

When $\Y$ has atoms, the stop-loss definition remains the primitive. Let
\[
q_\alpha=\inf\{t:F(t)\ge\alpha\},
\qquad
F(q_\alpha-)=\Pp(\Y<q_\alpha).
\]
The atom-split lower-tail expected shortfall is
\begin{equation}
\ES_{\alpha}
=\frac{1}{\alpha}\left(
\E[\Y\one\{\Y<q_\alpha\}]
+q_\alpha(\alpha-\Pp(\Y<q_\alpha))
\right).
\label{eq:discrete-es}
\end{equation}
The boundary atom is split so exactly $\alpha$ probability mass enters the
tail average.

\begin{proposition}[Discrete stop-loss value]
\label{prop:formal-discrete-stoploss}
If
\[
\Pp(\Y<q_\alpha)\le\alpha\le\Pp(\Y\le q_\alpha),
\]
then the optimized stop-loss value equals \cref{eq:discrete-es}.
\end{proposition}

\begin{twocolproof}
\twocolrow{
\[
\E[(q_\alpha-\Y)_+]
=\E[(q_\alpha-\Y)\one\{\Y<q_\alpha\}].
\]
}{
Only outcomes strictly below the boundary atom contribute to the lower-tail
shortfall.
}
\twocolrow{
\[
\Psi_\alpha(q_\alpha)
=q_\alpha-\frac{1}{\alpha}
\E[(q_\alpha-\Y)\one\{\Y<q_\alpha\}].
\]
}{
Substitute into the stop-loss objective.
}
\twocolrow{
\[
\Psi_\alpha(q_\alpha)
=q_\alpha-\frac{1}{\alpha}
\left(q_\alpha\Pp(\Y<q_\alpha)-\E[\Y\one\{\Y<q_\alpha\}]\right).
\]
}{
Expand the expectation.
}
\twocolrow{
\[
\Psi_\alpha(q_\alpha)
=\frac{1}{\alpha}
\left(\E[\Y\one\{\Y<q_\alpha\}]
+q_\alpha(\alpha-\Pp(\Y<q_\alpha))\right).
\]
}{
Rearrange. This is the atom-split formula.
}
\end{twocolproof}

For sorted observations $Y_{(1)}\le\cdots\le Y_{(n)}$ and
$k=\lfloor \alpha n\rfloor$, the sample atom-split lower-tail CVaR is
\begin{equation}
\widehat{\CVaR}_\alpha(\Y)
=\frac{1}{\alpha n}\left[
\sum_{i=1}^{k}Y_{(i)}+(\alpha n-k)Y_{(k+1)}
\right].
\label{eq:atom-sample}
\end{equation}
This is the full-oracle truth formula for tied rubric outcomes. The naive
average over all observations satisfying $Y_i\le q_\alpha$ overcounts the
boundary atom.

\subsection{\texorpdfstring{$\alpha=1$}{alpha=1} Collapse}
\label{app:alpha-one}

\begin{theorem}[$\alpha=1$ collapses to Mean-CJE]
\label{prop:formal-alpha-one}
If $\Y\in[0,1]$, then $\CVaR_1(\ptarget)=V(\ptarget)$. Taking any
$t\ge1$ and $h_t(\Sscore,\X)=t-f(\Sscore,\X)$, the CVaR shortfall residual
is the negative of the mean-CJE residual:
\[
\shortfall{t}{\Y}-h_t(\Sscore,\X)
=f(\Sscore,\X)-\Y.
\]
The tail-mass moment is identically zero when $t\ge1$.
\end{theorem}

\begin{twocolproof}
\twocolrow{
\[
\Psi_{1,\ptarget}(t)
=t-\E_\ptarget[\shortfall{t}{\Y}]
=\E_\ptarget[\min\{t,\Y\}].
\]
}{
For $\alpha=1$, $t-(t-\Y)_+=\min\{t,\Y\}$ pointwise.
}
\twocolrow{
\[
\E_\ptarget[\min\{t,\Y\}]\le\E_\ptarget[\Y],
\]
with equality for every $t\ge1$.
}{
The oracle outcome is bounded in $[0,1]$.
}
\twocolrow{
\[
\CVaR_1(\ptarget)=\E_\ptarget[\Y]=V(\ptarget).
\]
}{
Take the supremum over $t$.
}
\twocolrow{
\[
h_t(\Sscore,\X)
=\E_{\pzero}[t-\Y\mid\mathcal{G}]
=t-f(\Sscore,\X).
\]
}{
For $t\ge1$, the shortfall is $t-\Y$ and conditional expectation is linear.
}
\twocolrow{
\[
(t-\Y)-(t-f(\Sscore,\X))=f(\Sscore,\X)-\Y.
\]
}{
Substitute into the shortfall residual.
}
\twocolrow{
\[
\one\{\Y\le t\}-1=0.
\]
}{
If $t\ge1$ and $\Y\in[0,1]$, every observation lies below the threshold.
}
\end{twocolproof}

The identity is also an implementation invariant. At $\alpha=1$, the grid must
include a threshold above $\max_iY_i$; the fitted stop-loss calibrator must
equal $t-\hat f(\Sscore,\X)$; the CVaR point estimate, bootstrap draws,
calibration variance, and audit residuals must match the corresponding
Mean-CJE quantities up to numerical tolerance.
